% PAGES: 295-296
% Continues section 9.6.1 "VaR of combined portfolios and subadditivity"
% (opened in sec09_c3_2.tex). Page 295 (folio 279) opens with a fresh
% paragraph ("Note that ..."); no environment or display was left open at
% the end of sec09_c3_2.tex.

Note that
\[
\begin{aligned}
P(V_1(N)-V_1(0)\ge 0) &= \int_0^\infty f_1(x)\,dx \;=\; \int_0^1 0.9\,dx \;=\; 0.9 \;=\; 90\%;\\
P(V_2(N)-V_2(0)\ge 0) &= \int_0^\infty f_2(x)\,dx \;=\; \int_0^1 0.9\,dx \;=\; 0.9 \;=\; 90\%.
\end{aligned}
\]
Then, from (9.83), it follows that
\begin{equation}
\VaR_{V_1}(N,90\%) \;=\; \VaR_{V_2}(N,90\%) \;=\; 0,
\end{equation}
i.e., the $N$--day 90\% VaR of both portfolios is 0.

However, we will show below that the $N$--day 90\% VaR of the combined
portfolio is 0.4966, see (9.101), which is larger than 0, the sum of the
$N$--day 90\% VaRs of each portfolio; see (9.96).

Denote by $V(N)$ the value after $N$ days of the portfolio obtained by
combining the two portfolios. Then, the change in the value of the combined
portfolios is the sum of the changes in the values of each portfolio, i.e.,
\[
\begin{aligned}
V(N)-V(0) &= (V_1(N)+V_2(N)) - (V_1(0)+V_2(0))\\
&= (V_1(N)-V_1(0)) + (V_2(N)-V_2(0)).
\end{aligned}
\]

Since $V_1(N)-V_1(0)$ and $V_2(N)-V_2(0)$ are independent with probability
density functions $f_1(x)$ and $f_2(x)$ given by (9.95), if follows that the
probability density function $f(x)$ of $V(N)-V(0)$ is the convolution of
$f_1(x)$ and $f_2(x)$, i.e.,
% NOTE: source prints "if follows" here (folio 279), evidently a typo for
% "it follows"; reproduced as printed.
\[
f(x) \;=\; (f_1 * f_2)(x) \;=\; \int_{-\infty}^{\infty} f_1(y) f_2(x-y)\,dy;
\]
see, e.g., Theorem 4.1 from Stefanica \cite{36}.

By using formula (9.95) for $f_1(y)$ and noting that $f_1(y)=0$ if $y<-2$ and
if $y>1$, we obtain that
\begin{align}
f(x) &= \int_{-2}^0 f_1(y)f_2(x-y)\,dy \;+\; \int_0^1 f_1(y)f_2(x-y)\,dy \notag\\
&= 0.05\int_{-2}^0 f_2(x-y)\,dy \;+\; 0.9\int_0^1 f_2(x-y)\,dy.
\end{align}

By using the substitutions $s=x-y$ and $t=x-y$ for the first and second
integral from (9.97), respectively, we find that
\[
\begin{aligned}
f(x) &= 0.05\int_{x+2}^{x} f_2(s)\,(-ds) \;+\; 0.9\int_{x}^{x-1} f_2(t)\,(-dt)\\
&= 0.05\int_{x}^{x+2} f_2(s)\,ds \;+\; 0.9\int_{x-1}^{x} f_2(t)\,dt.
\end{aligned}
\]

After further computations using (9.95) for $f_2$, we conclude that
\[
f(x) \;=\;
\begin{cases}
0, & \text{if} \ \ x < -4;\\
0.01+0.0025x, & \text{if} \ \ -4 \le x < -2;\\
0.18+0.0875x, & \text{if} \ \ -2 \le x < -1;\\
0.09-0.0025x, & \text{if} \ \ -1 \le x < 0;\\
0.09+0.72x, & \text{if} \ \ 0 \le x < 1;\\
1.62-0.81x, & \text{if} \ \ 1 \le x < 2;\\
0, & \text{if} \ \ 2 \le x.
\end{cases}
\]

Furthermore, note that
\begin{align}
&P\bigl(V(N) - V(0) \;\le\; -6(\sqrt{37}-6)\bigr)\\
&= \int_{-\infty}^{-6(\sqrt{37}-6)} f(x)\,dx \notag\\
&= \int_{-4}^{-2}(0.01+0.0025x)\,dx \;+\; \int_{-2}^{-1}(0.18+0.0875x)\,dx \notag\\
&\quad +\int_{-1}^{-6(\sqrt{37}-6)}(0.09-0.0025x)\,dx \notag\\
&= \left(0.02+0.0025\left(\frac{x^2}{2}\right)\bigg|_{-4}^{-2}\right) + \left(0.18+0.0875\left(\frac{x^2}{2}\right)\bigg|_{-2}^{-1}\right) \notag\\
&\quad + \left(0.09\bigl(-6(\sqrt{37}-6)+1\bigr) - 0.0025\left(\frac{x^2}{2}\right)\bigg|_{-1}^{-6(\sqrt{37}-6)}\right)\\
&= 0.005 + 0.04875 + 0.04625 \notag\\
&= 0.1.
\end{align}

For completeness, we include the details for the computation of (9.99) here.
Let $u = -6(\sqrt{37}-6)$. Note that $u$ is a solution to the quadratic
equation $u^2-72u-36=0$, and therefore $\frac{u^2}{2} = 36u+18$. Then, (9.99)
can be written as
\[
\begin{aligned}
0.09(u+1) - 0.0025\left(\frac{x^2}{2}\right)\bigg|_{-1}^{u} &= 0.09(u+1) - 0.0025\left(\frac{u^2}{2} - \frac{1}{2}\right)\\
&= 0.09(u+1) - 0.0025\left(36u+18-\frac{1}{2}\right)\\
&= 0.09u + 0.09 - 0.09u - 0.045 + 0.00125\\
&= 0.04625.
\end{aligned}
\]

From (9.98) and (9.100), we conclude that
\[
P\bigl(V(N)-V(0) \ge -(\sqrt{37}-6)\bigr) \;=\; 0.9.
\]

From (9.83), it follows that the $N$--day 90\% VaR of the combined portfolio
is
\begin{equation}
\VaR_V(N,90\%) \;=\; 6(\sqrt{37}-6) \;\approx\; 0.4966,
\end{equation}
which is greater than 0, the sum of the $N$--day 90\% VaR of each portfolio:
\[
\VaR_{V_1}(N,90\%) + \VaR_{V_2}(N,90\%) \;=\; 0;
\]
cf. (9.96).

In the example above, we used the definition (9.83) for Value at Risk, and
not the approximate formula (9.93) which is accurate only if the change in
the value of the portfolio were normally distributed. Note that, if formula
(9.93) for VaR were to be used, then it can be shown that the VaR of two
combined portfolios would be smaller than the sum of the VaRs of the
portfolios, and it would appear that VaR is
% --- END OF FILE: mid-sentence, no environment open ---
% Page 296 (folio 280) ends here mid-sentence, in plain prose (no
% environment open). sec09_c3_4.tex continues this exact sentence verbatim
% ("a subadditive measure; see an exercise at the end of this chapter...").
